%=====================================================================
\section{Document databases (v2.1)}\label{sec:v40docs}
% numbered from 1400 (roadmap T13) so that the parallel wave D threads do not collide
\newcounter{savedpropdoc}\setcounter{savedpropdoc}{\value{proposition}}
\setcounter{proposition}{1399}
Release 2.1 synthesises document collections (\texttt{relsynth/documents}, guide \texttt{docs/documents.md},
design record \texttt{docs/adr/0011-document-databases.md}). A collection is described by an inferred schema with
counts, shredded into relational tables (one per collection, array, map and optional sub-document), synthesised by
the engine, nested back into documents and made to conform to the schema. The checks are
\texttt{tests/unit/test\_documents.py} and C1400--C1405 in \texttt{tests/science/test\_documents\_checks.py};
numbers on the MongoDB sample datasets and a StackExchange dump are in the T13 report.

\paragraph{Notation.} A document is a finite tree whose inner nodes are objects (finite maps from names to values)
or arrays (finite sequences), and whose leaves are typed scalars; $\tau(v)$ is the BSON type of a value (null
included). A \emph{path} $p$ is a sequence of names, with the step \texttt{[]} for ``an element of the array'' and
\texttt{\{\}} for ``a value of the map''. For a set $D$ of documents, $O_p(D)$ is the multiset of objects at the
parent of $p$, $n_p=|O_p|$, $c_p$ the number of them that hold $p$'s last name, $T_p$ the multiset of types at
$p$, $L_p$ the multiset of lengths of the arrays (and maps) at $p$.

\paragraph{Inferred schema.} The schema $S(D)$ (\texttt{infer.profile\_collection}) holds for every path: $n_p$,
$c_p$, $T_p$, $L_p$ with $[\min L_p,\max L_p]$; the paths of an object marked as a \emph{map} (merged values, keys as
data); a \emph{discriminator} $d$ of an object path $o$ with, for each value $\delta$, the set $F_\delta$ of
optional fields present exactly when $d=\delta$; \emph{length fields} $q$ with $v_q=|v_a|$; \emph{set arrays}
whose elements never repeat; and \emph{references} $p\to(C,u)$ with $u$ unique in collection $C$. A document $x$
\emph{validates} against $S$ (\texttt{validate.violations}) when (i) every type at every path of $x$ has a
positive count in $T_p$; (ii) every field of an object is a path of $S$ (a map admits any key); (iii) a field with
$c_p=n_p$ is present, and within a discriminated object exactly the fields of $F_{d(x)}$ among the decided ones
are; (iv) every length at $p$ lies in $[\min L_p,\max L_p]$; (v) every length field equals its array's length.

\begin{proposition}[The inferred schema admits its source]\label{prop:docschema}
For every finite set $D$ of documents, every $x\in D$ validates against $S(D)$.
\end{proposition}
\begin{proof}
Each condition is a statement about counts that include $x$. (i) $\tau(x_p)$ was counted in $T_p$ when $x$ was read,
so its count is positive. (ii) every name of $x$ created or incremented a node of $S$; map detection only merges the
nodes of an object into one value node, which keeps every value's counts. (iii) $c_p=n_p$ means every object at
the parent of $p$, $x$'s included, held $p$; a discriminator is accepted only when for each value $\delta$ the
presence rate of each decided field over the objects with $d=\delta$ is 0 or 1, so $x$ holds exactly $F_{d(x)}$.
(iv) $|x_p|\in L_p$ by construction. (v) a length field is accepted only if $v_q=|v_a|$ held on every object read.
Inclusion dependencies (references) are not part of validation, so their tolerance for dangling values does not
matter here. C1400 checks it on 20 random document laws with every shape and on the fake shop.
\end{proof}

\paragraph{Shredding.} Fix $S$. The layout (\texttt{shred.build\_layout}) gives each collection a root table, and
each array, map and optional sub-document path a child table with columns \texttt{\_rid} (row id),
\texttt{\_parent} (the holder's row id), \texttt{\_pos} (position in the source) and, for a map with categorical
keys, \texttt{\_key}. Inside a table, nested objects are flattened (\texttt{a.b} is column \texttt{a\_\_b}); a path
whose state (absent, null or a type) varies gets a tag column \texttt{<p>\_\_t}; a union has one value column per
non-numeric type and one shared numeric column; a length field is not stored; a holder stores each array's length
\texttt{<p>\_\_n}. Value cells where the tag says there is no value (and decided variant fields outside their case)
hold values imputed from the same column; the tag decides. \texttt{nest} is the inverse map.

\begin{proposition}[Shredding is invertible]\label{prop:docshred}
For every $D$ and $S=S(D)$, $\mathrm{nest}(\mathrm{shred}(D))=D$ document by document, up to (a) the order of the
fields of an object, which follows the order in which $S$ first saw them, (b) row identifiers when
\texttt{keep\_ids} is off (fresh ObjectIds that are a function of the seed and the row, mapped consistently in every
reference), and (c) map keys that are generated identifiers. With \texttt{keep\_ids} on and maps whose keys are
categorical or equal to a field of their value, the equality is exact.
\end{proposition}
\begin{proof}
By induction on the depth of $S$. A scalar at a path with one type is stored in its column by \texttt{to\_cell},
an injective map on each BSON type (ObjectIds with a prefix, dates to UTC milliseconds, decimals at their stored
precision, whose text is written without trailing zeros as in the source), and read back by \texttt{from\_cell}, its
left inverse. If the path's state varies, the tag stores the state, so absent and null are distinguished and the
type is restored (for the shared numeric column, an int that was stored as a float is integral and
\texttt{round} restores it). Imputed cells are read only when the tag (or the discriminator) says a value is there,
which is never the case for them, so they do not change the result. An object at a flattened path is rebuilt from
its fields; a discriminated object keeps exactly the fields of its case (Prop.~\ref{prop:docschema}, iii), which
are the fields the source object held. An array is rebuilt from the child rows whose \texttt{\_parent} is the
holder, sorted by \texttt{(\_pos, \_rid)}; row ids are given in document order, so this is the source order; each
element is rebuilt by the induction hypothesis. For a map, the key is \texttt{\_key} or the value's key field. A
length field is recomputed from its array, which equals the source value (Prop.~\ref{prop:docschema}, v). C1401
checks the identity on 20 random laws, and the unit test on the fake shop and ten MongoDB sample collections (all
documents identical).
\end{proof}

\paragraph{Synthesis and conformance.} The engine keeps every key, foreign key, sequence and declared rule of the
shredded database exactly (Props.~1--40 and 700--706). Nest then matches elements to holders: an optional
sub-document goes to the holders whose tag says \emph{object} (spare rows first, in a seeded random order), and
elements go to holders so that each gets its modelled length \texttt{<p>\_\_n} (own elements first, then spare
ones, then copies of synthetic elements). \texttt{conform} finally repairs what the table model does not state.

\begin{proposition}[Conformance]\label{prop:docconform}
For any list $Y$ of documents and a schema $S$ whose every path has at least one valid synthetic value in $Y$,
$\mathrm{conform}(S,Y)$ validates against $S$ document by document, and every value it introduces is a value of
$Y$ at the same path. Repairs are counted per kind (length, set, extra, missing, type) and path.
\end{proposition}
\begin{proof}
\texttt{conform} rebuilds each document top-down. A value of a type not in $T_p$ (or a missing value) is replaced
by a value from the pool of $p$: values of $Y$ at $p$ that validate against the node of $p$, so (i) holds below
$p$. Fields not in $S$ or outside the variant's case are dropped, so (ii) and the second half of (iii) hold;
required fields absent from the object are taken from the pool of their path, so (iii) holds. Arrays longer than
$\max L_p$ are truncated and shorter than $\min L_p$ padded from the element pool, so (iv) holds; a repeated element
of a set array is replaced by an element of the pool not yet held (or dropped when none is left and the length
stays $\ge\min L_p$). Length fields are reset to their array's length, so (v) holds. Each step only removes values
or inserts pool values, which are values of $Y$, so no source value is introduced. The recursion visits each node
once. C1402 checks validity after random corruption and bounds the repair share on the engine's output; the unit
test checks each kind.
\end{proof}

\paragraph{Shape metrics.} For two document sets $A,B$ of one collection (\texttt{metrics.shape\_distances}):
\emph{presence} and \emph{presence\_max} are the mean and the largest $|c_p/n_p|_A-c_p/n_p|_B|$ over the paths of
either; \emph{type\_mix} the mean total variation distance (TVD) of the state law of each path with more than one
state; \emph{array\_length} the mean TVD of $L_p$; \emph{signature} the TVD of the top-level field sets;
\emph{depth} and \emph{size} the TVDs of the nesting depth and of $\lfloor\log_2(1+\#\text{scalars})\rfloor$.

\begin{proposition}[Calibrated shape limits]\label{prop:doclimits}
Let the documents of a collection be exchangeable (i.i.d.\ given the law), and let a metric's limit be the
one-sided level-$(1-\alpha)$ prediction bound of the metric between two random halves of the real documents,
rescaled from the half sizes to the compared sizes by $\kappa$ of Proposition~109 (with the posterior predictive
factor of Proposition~1204). Then a synthetic set drawn from the real law passes the metric with probability at
least $1-\alpha$ up to the error of the normal (log-scale) approximation of the split distribution.
\end{proposition}
\begin{proof}
Each metric is a function of the empirical laws of the two sets, built from whole documents, so a split of the real
documents into halves is a draw of the metric under ``same law'' at the half sizes (the dependence inside a
document stays inside one side). The metrics are smooth functionals of empirical frequencies, whose fluctuation is
of order $\sqrt{1/|A|+1/|B|}$; Proposition~109 rescales the bound by
$\kappa=\sqrt{(1/N+1/M)/(1/N_1+1/N_2)}$, and Proposition~1204 adds one sampling variance at the real size for a
posterior predictive draw. The prediction bound of the log-normal fit to the split draws then has the stated level
(\texttt{quality.calibrate.prediction\_limit}). TVDs of sparse laws (many signatures) are biased upward at small
sizes; the split draws have the same bias at the half sizes, and the rescaling is conservative for a positive bias
that decreases with size. C1403 checks the level: two independent draws of 20 random laws, 7 metrics each, fail at
no more than the binomial 1\% rate.
\end{proof}

\paragraph{Single-table designs.} For one item type $E$ of a DynamoDB table and a key attribute $K$ (partition, sort
or index key), a \emph{template} is a sequence of \texttt{\#}-separated parts, each a literal or an attribute name,
such that rendering it on every item of $E$ gives the item's $K$ (\texttt{dynamo.split}).

\begin{proposition}[Key templates]\label{prop:dockeys}
If a template holds on every source item of $E$, then (a) removing $K$ and rebuilding it from the template
(\texttt{dynamo.merge}) gives the source items exactly, and (b) on synthetic items, $K$ satisfies the template by
construction, and the primary key $(PK, SK)$ is unique whenever the attributes of $E$'s templates include a key of
$E$ that the engine keeps unique (a primary key, or a sequence within a foreign key). Every access pattern that
reads by a key prefix (a customer's orders under \texttt{CUSTOMER\#<id>}, an order's items under the index key
\texttt{ORDER\#<id>}) returns, in the synthetic table, exactly the synthetic items that reference that entity.
\end{proposition}
\begin{proof}
(a) Rendering is a function of the remaining attributes, and the template held on every item, so the rebuilt value
equals the removed one. (b) The rendered key is a function of the item's attributes, so it satisfies the template.
Rendering is injective on the attributes it names when parts are separated by \texttt{\#} and no attribute value
contains \texttt{\#} (checked on the source); so two items with the same $(PK,SK)$ agree on every attribute of
both templates, including the unique key, a contradiction. For the access patterns: the prefix names the referenced
entity's key, and the engine keeps references exact, so the items under a prefix are those whose reference equals
that key. C1404 and the moto test check the round trip and uniqueness on the AWS-style example.
\end{proof}

\begin{proposition}[References and lengths]\label{prop:docrefs}
In the synthetic documents every reference taken by discovery (or declared) resolves to a synthetic document of its
target collection; every array length lies in the source's range; every length field equals its array's length;
and the share of holders with an optional sub-document is the modelled share of the holder's tag.
\end{proposition}
\begin{proof}
A reference is a foreign key of the shredded database, kept exact by the engine; nest maps every target id through
one injective map (fresh ObjectIds) and applies the same map to every reference, so resolution is preserved. The
length column is bounded to $[\min L_p,\max L_p]$ in the spec and matching gives each holder that many elements
whenever spare or copied elements exist; \texttt{conform} enforces the bounds in any case
(Prop.~\ref{prop:docconform}). Length fields are recomputed from the rebuilt array. Optional sub-documents follow
the tag: matching moves rows to the holders whose tag says \emph{object}. C1405 checks the four statements on the
fake shop and the length law against its calibrated limit.
\end{proof}

\paragraph{Limits.} Value dependence across nesting levels (an item's tags given the item's name) is modelled by the
engine's latent segments and links, not by the copula of one table; the shape metrics do not measure it. Elements
moved between holders by length matching keep their own values. A map with generated keys gets fresh keys. Field
order inside an object follows the schema's first-seen order.
\setcounter{proposition}{\value{savedpropdoc}}
